%% Provisional abstract (2026-09-10). [E2]/[E3] clauses are filled from the
%% experiment artifacts only; the registered abstract (Sep 18) must match.
Serial spectral--spatial models compress each pixel's spectrum with a small
encoder and classify with a wide spatial network that reads the encoded
neighbourhood. When the neighbourhood predicts the label, does joint
training starve the encoder? We prove, in a linear-encoder cross-entropy
model with $M$ replicated contextual readouts, that at a matched fitting
margin the spectral coefficient moves by at most a fraction
$1/(1+Mv_0^2)$ of the update a spectral-only model needs, that the fitted
model then fails under contextual reversal when its initial spectral
coefficient is below half the margin, with expected error tending to
$\Phi(m/2\sigma)$ over isotropic initialization, and that a
$M^{-1/2}$-normalized readout removes the width dependence without
removing the failure. The relative training speed of the contextual
pathway, not width as such, is the competition parameter. We test the
prediction with interventions on that speed in a ReLU convolutional model
with two cues of equal information but unequal initial accessibility [E2:
verdict], and with measured tissue spectra in constructed neighbourhoods
[E3: verdict], each with a control that removes the cause. On a production
model, scalar curvature does not identify the mechanism and an earlier
apparent freezing advantage was protocol-confounded; we recommend no
mitigation.
